% =========================================================================
% Secao 1 - Introducao
% =========================================================================

\section{Introdução}
\label{sec:introducao}

O Brasil aboliu a obrigatoriedade do registro de matrícula em séries
específicas em 1971 e, desde então, convive com um problema mensurável
apenas indiretamente. O fluxo escolar, isto é, a passagem de coortes de
estudantes por séries, modalidades e redes ao longo do tempo, depende de
um sistema de registros administrativos que captura imperfeitamente
trajetórias individuais. Por décadas, o Censo Escolar do Instituto
Nacional de Estudos e Pesquisas Educacionais Anísio Teixeira (INEP)
divulgou taxas de rendimento escolar, isto é, aprovação, reprovação e
abandono, baseadas em estoques agregados de matrícula, sem
acompanhamento longitudinal do mesmo aluno. Estas taxas medem a
situação do aluno ao final do ano letivo e são apuradas anualmente.

A introdução do identificador único do aluno via CPF em 2007 permitiu
ao INEP construir um conjunto distinto e complementar de medidas. Trata-se
das taxas de transição, também chamadas de indicadores de fluxo, que
medem a passagem do aluno entre dois anos letivos consecutivos em
quatro categorias, a saber, promoção, repetência, migração para EJA e
evasão. As duas famílias de indicadores não devem ser confundidas. A
medida de rendimento captura o que aconteceu dentro de um ano letivo,
enquanto a medida de fluxo captura a passagem entre anos. A última
transição divulgada pelo INEP refere-se aos anos 2021 e 2022, publicada
em 2025, e o ritmo das publicações posteriores é desconhecido. O INEP
não estabelece, em sua página oficial, calendário público para
atualização desse indicador.\footnote{A nota técnica original do INEP
cobre o período 2007--2016 \citep{INEP2017_fluxo}. A série foi atualizada
até a transição 2021--2022 \citep{INEP_transicao}. Em junho de 2026,
não há atualizações para transições posteriores.}

A distinção entre rendimento e fluxo importa porque a repetência
efetiva e a evasão genuína só podem ser identificadas longitudinalmente.
O aluno reprovado em $t$ pode repetir a mesma série em $t+1$, mas
também pode mudar de escola, migrar para EJA ou evadir, e os indicadores
de rendimento sozinhos não distinguem entre essas trajetórias. Quanto às
taxas de rendimento propriamente ditas, a última divulgação refere-se ao
ano letivo 2024, publicada em agosto de 2025, com defasagem típica de
oito a dez meses.

Este artigo propõe construir os indicadores de fluxo escolar utilizando
uma fonte de dados que oferece vantagens estruturais sobre o Censo
Escolar, a saber, o painel longitudinal da Pesquisa Nacional por
Amostra de Domicílios Contínua (PNADC) do Instituto Brasileiro de
Geografia e Estatística (IBGE). Desde sua introdução em 2012, a PNADC
adotou um \emph{design} amostral em que cada domicílio é visitado em
cinco trimestres consecutivos. A liberação recente dos identificadores
longitudinais individuais por parte do IBGE permite, pela primeira vez,
observar diretamente o status de matrícula do mesmo aluno ao longo de
até quinze meses, janela suficiente para distinguir transições
intra-ano, associadas ao abandono, de transições entre anos, que
correspondem à evasão, à promoção e à repetência.

Trabalhos prévios na literatura brasileira já enfrentaram, de outras
formas, a tarefa de medir fluxo escolar a partir de pesquisa
domiciliar. O esforço mais conhecido nessa direção é o do grupo
PROFLUXO, sigla do Programa de Fluxo Escolar, proposto por
\cite{Fletcher1985_modelo} e desenvolvido na Fundação Cesgranrio por
Sergio Costa Ribeiro, Ruben Klein e colaboradores. O PROFLUXO, em
seus diversos desdobramentos \citep{Fletcher1985_repetencia,
Fletcher1985_modelo, FletcherRibeiro1989_profluxo,
Ribeiro1991_pedagogia, KleinRibeiro1991_censo, Klein2006_educacao},
estabeleceu que as estimativas de fluxo escolar derivadas diretamente
do Censo Escolar subestimavam sistematicamente a repetência e propôs
uma alternativa baseada em modelo de fluxo de estado estacionário
aplicado à matriz idade-série de um único corte da PNAD. A
contribuição histórica do grupo é amplamente reconhecida, em
particular o achado emblemático de \cite{Ribeiro1991_pedagogia} sobre
a pedagogia da repetência no Brasil. Esta nota não é uma extensão do
PROFLUXO. Não há sobreposição entre os autores, e os instrumentos
metodológicos são distintos. Onde o PROFLUXO usava um modelo de
estado estacionário para inferir fluxos não observáveis a partir de
um único corte, esta nota observa diretamente as transições entre
dois pontos no tempo no painel longitudinal da PNADC. A menção ao
PROFLUXO se justifica pelo lugar do grupo na história das
estatísticas educacionais brasileiras, mas o arcabouço técnico desta
nota é completamente diferente do que aquele grupo desenvolveu.

A contribuição deste trabalho é tríplice. Em primeiro lugar, oferece-se
uma reformulação metodológica dos cinco indicadores de fluxo escolar,
a saber, abandono intra-ano e evasão, promoção, repetência e
não-progressão entre anos, construídos a partir do painel longitudinal
da PNADC entre 2012 e 2024, cobrindo 52 trimestres consecutivos. Em
segundo lugar, exploramos sistematicamente a heterogeneidade
socioeconômica das estimativas, isto é, por quintil de renda
domiciliar per capita, perfil de elegibilidade ao Bolsa Família,
raça/cor, sexo, defasagem idade-série, rede e macrorregião, demonstrando
dimensões de desigualdade que os agregados do INEP não revelam por
construção. Em terceiro lugar, comparamos as estimativas com as do
INEP em todos os anos em que esses são publicados, decompondo as
diferenças em cinco fontes, a saber, retorno, universo, amostragem,
sub-cobertura e medida residual.

A análise demonstra que a PNADC oferece quatro vantagens sobre as
estatísticas oficiais do INEP, todas mensuráveis. A primeira vantagem
diz respeito à tempestividade, pois a PNADC é divulgada trimestralmente,
enquanto a última atualização dos Indicadores de Trajetória do INEP,
relativa à transição 2021--2022, foi publicada em julho de 2025 e o
INEP não estabelece calendário público para atualizações posteriores.
A segunda vantagem é a desagregação socioeconômica, pois a PNADC tem
variáveis individuais de renda, raça e Bolsa Família que o Censo
Escolar não possui. A terceira vantagem reside na independência da
fonte, pois pesquisa domiciliar e registro administrativo são
metodologicamente distintos e validam-se mutuamente. Por fim, a
quarta vantagem é a captura de retorno, pois a PNADC acompanha o
indivíduo no domicílio mesmo quando ele muda de rede, unidade
federativa ou modalidade, enquanto o Censo perde alunos que migram
para escolas mal cobertas pelo registro.

O restante da nota organiza-se da seguinte forma. A
Seção~\ref{sec:tres_eras} situa os trabalhos prévios sobre fluxo
escolar no Brasil e a divulgação oficial dos indicadores. A Seção~\ref{sec:dados_metodo} descreve os microdados
da PNADC, o esquema amostral rotativo, a construção do painel
longitudinal e as definições operacionais dos cinco indicadores. As
Seções~\ref{sec:agregados} e~\ref{sec:heterogeneidade} apresentam,
respectivamente, os resultados agregados nacionais e a heterogeneidade
socioeconômica. A Seção~\ref{sec:inep} confronta as estimativas com as
do INEP e decompõe as diferenças. A Seção~\ref{sec:robustez} reporta
testes de robustez, com destaque para o atrito do painel, a regra de
seleção de trimestre de referência e a reponderação. A
Seção~\ref{sec:conclusao} sintetiza as quatro vantagens, discute
limitações e propõe uma agenda.
